\chapter{Rotating TFT Black Holes and the Kerr Correspondence Limit}
\label{app:kerr_limit}
\addtocontents{toc}{\protect\setcounter{tocdepth}{0}}

\noindent\textbf{Subject:} Torsional Soft-Floor Shear, Rotating Phase Boundaries, and the Exterior Correspondence Conditions for Kerr-like Observables.

The live-boundary formulation developed in Appendix Q assigns black-hole angular momentum to torsional motion of the FCC/SC phase boundary and to shear strain in the surrounding Soft Floor. That construction provides a mechanical carrier for spin. It does not, by itself, determine the full exterior field of a rotating black hole.

The purpose of this appendix is to define the exterior correspondence conditions under which the TFT continuum description can reproduce the successful external observables associated with Kerr black holes while retaining a different interior ontology. The Kerr solution is treated here as the reference exterior model for an uncharged rotating black hole. The corresponding TFT problem is to derive the external angular-momentum field from boundary torsion and Soft-Floor shear, rather than to postulate it geometrically.

The required correspondence is
\begin{equation}
    \text{Kerr angular momentum} \quad\longrightarrow\quad \text{FCC/SC boundary torsion} + \text{external Soft-Floor shear}.
\end{equation}

The analysis below is therefore a matching program. It identifies the external-field structure, boundary constraints, and observational signatures that a TFT Kerr limit must reproduce.

\section{Kerr as the Rotating Black-Hole Target}
A rotating black hole is characterized macroscopically by mass and axial angular momentum. In the TFT boundary description, the macroscopic angular momentum of the rotating black-hole system is denoted by
\begin{equation}
    \mathcal{J}_{\mathrm{BH}}.
\end{equation}
The corresponding dimensionless spin parameter is
\begin{equation}
    a_\ast = \frac{c_{\mathrm{E}}\mathcal{J}_{\mathrm{BH}}}{GM^2}.
\end{equation}

A Kerr-like TFT limit must reproduce more than the algebraic existence of $a_\ast$. The rotating exterior must also account for the principal spin-dependent structures associated with Kerr black holes:
\begin{itemize}
    \item frame dragging,
    \item spin-dependent horizon structure,
    \item ergosphere structure,
    \item spin-dependent ISCO behavior,
    \item photon-ring and shadow deformation,
    \item spin-dependent ringdown modes,
    \item restricted external multipole structure.
\end{itemize}

The bound $a_\ast \leq 1$ is only a first consistency condition. The central Kerr-limit problem is the derivation of the external torsional field generated by $\mathcal{J}_{\mathrm{BH}}$.

\section{Spin Normalization from the Live-Boundary Model}
The live-boundary appendix gives the surface-shell normalization for black-hole spin. In that model, axial angular momentum is carried by the FCC/SC boundary and by the external torsional response of the Soft Floor.

The imported relation is
\begin{equation}
    \mathcal{J}_{\mathrm{BH}} = I_{\mathrm{surf}} \frac{v_{\mathrm{tor}}}{R_{\mathrm{QG}}},
\end{equation}
where $I_{\mathrm{surf}}$ is the effective surface moment of inertia, $v_{\mathrm{tor}}$ is the torsional boundary speed, and $R_{\mathrm{QG}}$ is the TFT saturation radius.

Using
\begin{equation}
    R_{\mathrm{QG}} = \frac{GM}{c_{\mathrm{E}}^2},
\end{equation}
the spin parameter becomes
\begin{equation}
    a_\ast = \left( \frac{I_{\mathrm{surf}}}{M R_{\mathrm{QG}}^2} \right) \frac{v_{\mathrm{tor}}}{c_{\mathrm{E}}}.
\end{equation}

For a boundary-confined surface current,
\begin{equation}
    \frac{I_{\mathrm{surf}}}{M R_{\mathrm{QG}}^2} \leq 1,
\end{equation}
and the Soft-Floor kinematic capacity gives
\begin{equation}
    v_{\mathrm{tor}} \leq c_{\mathrm{E}}.
\end{equation}
Thus,
\begin{equation}
    a_\ast \leq 1.
\end{equation}

This result fixes the spin normalization used below. The remaining Kerr-limit question is whether this boundary angular momentum sources an external torsional Soft-Floor field with the required Kerr-like structure.

\section{Axisymmetric Torsional Soft-Floor Field}

A rotating TFT black hole sources an azimuthal shear response in the surrounding FCC Soft Floor. This response is not ordinary material rotation of the SC saturated interior. The saturated SC domain supplies the phase-locked internal boundary condition, while the live FCC/SC boundary imposes torsional loading on the external Soft Floor.

Let the spin axis define the \(z\)-axis and let
\begin{equation}
    r_\perp = r\sin\theta
\end{equation}
denote the cylindrical distance from that axis. In the stationary axisymmetric exterior, the leading torsional response may be written as
\begin{equation}
    v_{\phi}^{(\mathrm{SF})}(r,\theta)
    =
    r_\perp\,\omega_{\mathrm{drag}}(r)
    =
    r\sin\theta\,\omega_{\mathrm{drag}}(r),
\end{equation}
where \(\omega_{\mathrm{drag}}(r)\) is the local torsional dragging rate of the Soft Floor.

Neighboring spherical shells are torsionally sheared when their dragging rates differ. The corresponding radial torsional shear intensity is
\begin{equation}
    \gamma_{r\phi}^{(\mathrm{tor})}
    =
    r\sin\theta
    \frac{d\omega_{\mathrm{drag}}}{dr}.
\end{equation}

Let \(\tau_{\mathrm{tor}}(r)\) denote the effective torsional stiffness of the external Soft Floor. (Here \(\tau_{\mathrm{tor}}(r)\) denotes the effective three-dimensional torsional stiffness density of the strained FCC Soft Floor in the azimuthal shear channel. It is not introduced as a new fundamental constant; in the weak exterior it approaches the baseline Soft-Floor torsional stiffness, while near the saturated boundary it may acquire strain dependence through \(\alpha_G(r)\).)

Then,tThe exterior torsional strain-energy cost is proportional to
\begin{equation}
    \mathcal{E}_{\mathrm{tor}}
    \propto
    \int_{R_{\mathrm{QG}}}^{\infty}
    \int
    \tau_{\mathrm{tor}}(r)
    \left[
    r\sin\theta
    \frac{d\omega_{\mathrm{drag}}}{dr}
    \right]^2
    r^2\sin\theta
    \,d\theta\,d\phi\,dr .
\end{equation}

Using
\begin{equation}
    \int_0^{2\pi}
    \int_0^\pi
    \sin^3\theta
    \,d\theta\,d\phi
    =
    \frac{8\pi}{3},
\end{equation}
this becomes
\begin{equation}
    \mathcal{E}_{\mathrm{tor}}
    \propto
    \frac{8\pi}{3}
    \int_{R_{\mathrm{QG}}}^{\infty}
    \tau_{\mathrm{tor}}(r)
    r^4
    \left(
    \frac{d\omega_{\mathrm{drag}}}{dr}
    \right)^2
    dr .
\end{equation}

The stationary exterior torsional field is obtained by minimizing this functional subject to the boundary rotation at \(R_{\mathrm{QG}}\). The Euler--Lagrange equation is
\begin{equation}
    \frac{d}{dr}
    \left[
    \tau_{\mathrm{tor}}(r)
    r^4
    \frac{d\omega_{\mathrm{drag}}}{dr}
    \right]
    =
    0 .
\end{equation}

This is the leading axisymmetric torsional Soft-Floor field equation.

The solution below uses the strong-coupling boundary condition in which the external FCC Soft Floor is torsionally locked to the live FCC/SC boundary at \(R_{\mathrm{QG}}\). Thus,
\begin{equation}
    \omega_{\mathrm{drag}}(R_{\mathrm{QG}})
    =
    \omega_{\partial},
\end{equation}
where
\begin{equation}
    \omega_{\partial}
    =
    \frac{v_{\mathrm{tor}}}{R_{\mathrm{QG}}}.
\end{equation}

This is the no-slip torsional limit of the boundary coupling. A more general formulation would prescribe the transmitted boundary torque rather than the boundary angular velocity. In that case, the conjugate exterior torsional loading is proportional to
\begin{equation}
    \left.
    \tau_{\mathrm{tor}}(r)
    r^4
    \frac{d\omega_{\mathrm{drag}}}{dr}
    \right|_{r=R_{\mathrm{QG}}},
\end{equation}
which is the radial torsional flux appearing inside the field equation. The present appendix uses the strong-coupling form because it gives the clean leading correspondence limit between the rotating boundary and the exterior Soft-Floor shear field.

In the weak exterior, where the torsional stiffness approaches its baseline value,
\begin{equation}
    \tau_{\mathrm{tor}}(r)\rightarrow \tau_{\mathrm{tor},0},
\end{equation}
the equation reduces to
\begin{equation}
    \frac{d}{dr}
    \left[
    r^4
    \frac{d\omega_{\mathrm{drag}}}{dr}
    \right]
    =
    0 .
\end{equation}

Imposing
\begin{equation}
    \omega_{\mathrm{drag}}(\infty)=0
\end{equation}
and
\begin{equation}
    \omega_{\mathrm{drag}}(R_{\mathrm{QG}})
    =
    \omega_{\partial}
    =
    \frac{v_{\mathrm{tor}}}{R_{\mathrm{QG}}},
\end{equation}
gives
\begin{equation}
    \omega_{\mathrm{drag}}(r)
    =
    \omega_{\partial}
    \left(
    \frac{R_{\mathrm{QG}}}{r}
    \right)^3 .
\end{equation}

Therefore,
\begin{equation}
    \omega_{\mathrm{drag}}(r)
    =
    \frac{v_{\mathrm{tor}}}{R_{\mathrm{QG}}}
    \left(
    \frac{R_{\mathrm{QG}}}{r}
    \right)^3 ,
\end{equation}
and the corresponding tangential Soft-Floor shear speed is
\begin{equation}
    v_{\phi}^{(\mathrm{SF})}(r,\theta)
    =
    v_{\mathrm{tor}}
    \left(
    \frac{R_{\mathrm{QG}}}{r}
    \right)^2
    \sin\theta .
\end{equation}

Using the live-boundary spin relation
\begin{equation}
    \mathcal{J}_{\mathrm{BH}}
    =
    I_{\mathrm{surf}}
    \frac{v_{\mathrm{tor}}}{R_{\mathrm{QG}}},
\end{equation}
the exterior dragging rate can also be written as
\begin{equation}
    \omega_{\mathrm{drag}}(r)
    =
    \frac{\mathcal{J}_{\mathrm{BH}}}{I_{\mathrm{surf}}}
    \left(
    \frac{R_{\mathrm{QG}}}{r}
    \right)^3 .
\end{equation}

Since
\begin{equation}
    R_{\mathrm{QG}}
    =
    \frac{GM}{c_{\mathrm{E}}^2},
\end{equation}
this becomes
\begin{equation}
    \omega_{\mathrm{drag}}(r)
    =
    \left(
    \frac{M R_{\mathrm{QG}}^2}{I_{\mathrm{surf}}}
    \right)
    \frac{G\mathcal{J}_{\mathrm{BH}}}{c_{\mathrm{E}}^2 r^3}.
\end{equation}

Thus, the TFT torsional Soft-Floor field naturally recovers the required far-field angular-momentum scaling,
\begin{equation}
    \omega_{\mathrm{drag}}(r)
    \propto
    \frac{G\mathcal{J}_{\mathrm{BH}}}{c_{\mathrm{E}}^2 r^3}.
\end{equation}

In terms of the dimensionless spin parameter,
\begin{equation}
    a_\ast
    =
    \frac{c_{\mathrm{E}}\mathcal{J}_{\mathrm{BH}}}{GM^2},
\end{equation}
the tangential Soft-Floor shear speed is
\begin{equation}
    v_{\phi}^{(\mathrm{SF})}(r,\theta)
    =
    c_{\mathrm{E}}a_\ast
    \left(
    \frac{M R_{\mathrm{QG}}^2}{I_{\mathrm{surf}}}
    \right)
    \left(
    \frac{R_{\mathrm{QG}}}{r}
    \right)^2
    \sin\theta .
\end{equation}

More generally, if the torsional stiffness varies near the saturated boundary, the solution of the field equation is
\begin{equation}
    \omega_{\mathrm{drag}}(r)
    =
    \omega_{\partial}
    \frac{
    \displaystyle
    \int_r^\infty
    \frac{ds}{\tau_{\mathrm{tor}}(s)s^4}
    }{
    \displaystyle
    \int_{R_{\mathrm{QG}}}^{\infty}
    \frac{ds}{\tau_{\mathrm{tor}}(s)s^4}
    } .
\end{equation}

The inverse-cubic dragging law is therefore the constant-stiffness exterior limit of the more general torsional Soft-Floor response.

\section{Effective Rotating Optical Geometry}

The torsional Soft-Floor field derived above determines how the exterior vacuum is mechanically dragged by a rotating FCC/SC boundary. To connect this field to observable propagation, the static TFT optical law must be extended to a rotating medium.

In the nonrotating exterior, the local propagation speed of a transverse excitation is
\begin{equation}
    c_P(r)
    =
    c_{\mathrm{E}}
    \left(1-\alpha_G(r)\right)^2,
\end{equation}
with
\begin{equation}
    \alpha_G(r)
    =
    \frac{GM}{r c_{\mathrm{E}}^2}
    =
    \frac{R_{\mathrm{QG}}}{r} .
\end{equation}
Thus,
\begin{equation}
    c_P(r)
    =
    c_{\mathrm{E}}
    \left(
    1-
    \frac{R_{\mathrm{QG}}}{r}
    \right)^2 .
\end{equation}

In the rotating exterior, the Soft Floor also carries the azimuthal drift
\begin{equation}
    v_{\phi}^{(\mathrm{SF})}(r,\theta)
    =
    r\sin\theta\,\omega_{\mathrm{drag}}(r).
\end{equation}
The local propagation problem therefore becomes a moving-medium problem: waves propagate through a radially strained medium that is also torsionally dragged in the azimuthal direction.

Let a high-frequency excitation have phase
\begin{equation}
    \Theta_{\mathrm{ray}}
    =
    -\omega t
    +
    m_{\phi}\phi
    +
    \Theta_r(r)
    +
    \Theta_{\theta}(\theta),
\end{equation}
where \(\omega\) is the conserved frequency and \(m_{\phi}\) is the azimuthal mode number. The local wavevector components are
\begin{equation}
    k_r
    =
    \frac{\partial \Theta_r}{\partial r},
    \qquad
    k_{\theta}
    =
    \frac{\partial \Theta_{\theta}}{\partial \theta},
    \qquad
    k_{\phi}
    =
    \frac{m_{\phi}}{r\sin\theta} .
\end{equation}

The frequency measured in the locally dragged Soft-Floor frame is shifted by the azimuthal drift:
\begin{equation}
    \omega_{\mathrm{loc}}
    =
    \omega
    -
    \mathbf{k}\cdot\mathbf{v}_{\mathrm{SF}} .
\end{equation}
Using
\begin{equation}
    \mathbf{k}\cdot\mathbf{v}_{\mathrm{SF}}
    =
    k_{\phi}v_{\phi}^{(\mathrm{SF})}
    =
    \frac{m_{\phi}}{r\sin\theta}
    \left(r\sin\theta\,\omega_{\mathrm{drag}}\right),
\end{equation}
gives
\begin{equation}
    \omega_{\mathrm{loc}}
    =
    \omega
    -
    m_{\phi}\omega_{\mathrm{drag}}(r).
\end{equation}

The local Soft-Floor dispersion relation is therefore
\begin{equation}
    \left[
    \omega
    -
    m_{\phi}\omega_{\mathrm{drag}}(r)
    \right]^2
    =
    c_P^2(r)
    \left[
    k_r^2
    +
    \frac{k_{\theta}^2}{r^2}
    +
    \frac{m_{\phi}^2}{r^2\sin^2\theta}
    \right].
\end{equation}

This is the leading eikonal propagation law for the rotating TFT exterior. The same torsional field that carries black-hole angular momentum now shifts the local wave frequency by \(m_{\phi}\omega_{\mathrm{drag}}\), thereby producing the TFT frame-dragging mechanism. Prograde modes \((m_{\phi}\omega_{\mathrm{drag}}>0)\) and retrograde modes \((m_{\phi}\omega_{\mathrm{drag}}<0)\) therefore propagate through different effective strain environments.

The same result may be packaged as an effective rotating optical line element. Let the spatial optical factor remain the static TFT strain factor and replace the azimuthal differential by the locally dragged differential \(d\phi-\omega_{\mathrm{drag}}dt\). The resulting leading stationary axisymmetric optical form is
\begin{equation}
\begin{aligned}
    ds_{\tau,\mathrm{rot}}^2
    ={}&
    -c_{\mathrm{E}}^2
    \left(1-\alpha_G\right)^2dt^2
    +
    \left(1-\alpha_G\right)^{-2}
    \Big[
    dr^2
    +
    r^2d\theta^2
    \\
    &
    +
    r^2\sin^2\theta
    \left(
    d\phi
    -
    \omega_{\mathrm{drag}}(r)dt
    \right)^2
    \Big].
\end{aligned}
\end{equation}

Expanding the azimuthal term gives the effective off-diagonal component
\begin{equation}
    g_{t\phi}^{(\tau)}
    =
    -
    \left(1-\alpha_G\right)^{-2}
    r^2\sin^2\theta\,
    \omega_{\mathrm{drag}}(r).
\end{equation}
Thus, the off-diagonal frame-dragging structure is not postulated geometrically. It is induced by the torsional motion of the external Soft Floor.

Using the constant-stiffness exterior solution,
\begin{equation}
    \omega_{\mathrm{drag}}(r)
    =
    \omega_{\partial}
    \left(
    \frac{R_{\mathrm{QG}}}{r}
    \right)^3,
\end{equation}
the off-diagonal exterior component scales as
\begin{equation}
    g_{t\phi}^{(\tau)}
    \propto
    -
    \left(1-\alpha_G\right)^{-2}
    \frac{R_{\mathrm{QG}}^3\omega_{\partial}}{r}
    \sin^2\theta .
\end{equation}
Equivalently, using the angular-momentum form,
\begin{equation}
    g_{t\phi}^{(\tau)}
    \propto
    -
    \left(1-\alpha_G\right)^{-2}
    \left(
    \frac{M R_{\mathrm{QG}}^2}{I_{\mathrm{surf}}}
    \right)
    \frac{G\mathcal{J}_{\mathrm{BH}}}{c_{\mathrm{E}}^2 r}
    \sin^2\theta .
\end{equation}

The important result is the radial structure: the angular dragging rate falls as \(r^{-3}\), while the off-diagonal optical coupling falls as \(r^{-1}\) in the far exterior. This is the required leading behavior for a rotating angular-momentum field in a Kerr-correspondence limit. The remaining difference between this TFT optical form and the exact Kerr metric is a strong-field normalization problem, not the absence of a mechanical frame-dragging carrier.

\section{Torsional Capacity Surface and Ergosphere Correspondence}

The rotating optical form above gives a mechanical criterion for the TFT analogue of an ergosphere. In a stationary axisymmetric exterior, a locally nonrotating state becomes mechanically unavailable when the azimuthal Soft-Floor drift equals the local transverse propagation capacity.

The torsional capacity condition is therefore
\begin{equation}
    \left|
    v_{\phi}^{(\mathrm{SF})}(r,\theta)
    \right|
    =
    c_P(r).
\end{equation}
Using
\begin{equation}
    v_{\phi}^{(\mathrm{SF})}(r,\theta)
    =
    v_{\mathrm{tor}}
    \left(
    \frac{R_{\mathrm{QG}}}{r}
    \right)^2
    \sin\theta
\end{equation}
and
\begin{equation}
    c_P(r)
    =
    c_{\mathrm{E}}
    \left(
    1-
    \frac{R_{\mathrm{QG}}}{r}
    \right)^2,
\end{equation}
the capacity surface is defined by
\begin{equation}
    v_{\mathrm{tor}}
    \left(
    \frac{R_{\mathrm{QG}}}{r}
    \right)^2
    \sin\theta
    =
    c_{\mathrm{E}}
    \left(
    1-
    \frac{R_{\mathrm{QG}}}{r}
    \right)^2 .
\end{equation}

Introduce the dimensionless radius
\begin{equation}
    x
    =
    \frac{r}{R_{\mathrm{QG}}}
\end{equation}
and the boundary torsional fraction
\begin{equation}
    u_{\mathrm{tor}}
    =
    \frac{v_{\mathrm{tor}}}{c_{\mathrm{E}}} .
\end{equation}
The capacity condition becomes
\begin{equation}
    u_{\mathrm{tor}}
    \frac{\sin\theta}{x^2}
    =
    \left(
    1-
    \frac{1}{x}
    \right)^2
    =
    \frac{(x-1)^2}{x^2} .
\end{equation}
Multiplying by \(x^2\) gives
\begin{equation}
    (x-1)^2
    =
    u_{\mathrm{tor}}\sin\theta .
\end{equation}
The outer torsional capacity surface is therefore
\begin{equation}
    r_{\mathrm{cap}}(\theta)
    =
    R_{\mathrm{QG}}
    \left[
    1
    +
    \sqrt{
    u_{\mathrm{tor}}\sin\theta
    }
    \right].
\end{equation}

In this formulation, the ergosphere correspondence is not introduced as a metaphor. It is the region in which the Soft Floor is dragged so strongly that nonrotating propagation is no longer kinematically available. On the spin axis, \(\sin\theta=0\), and the torsional capacity surface collapses to the saturation boundary:
\begin{equation}
    r_{\mathrm{cap}}(0)
    =
    R_{\mathrm{QG}} .
\end{equation}
At the equator,
\begin{equation}
    r_{\mathrm{cap}}
    \left(\frac{\pi}{2}\right)
    =
    R_{\mathrm{QG}}
    \left(
    1+
    \sqrt{u_{\mathrm{tor}}}
    \right).
\end{equation}
For the limiting boundary speed \(v_{\mathrm{tor}}\rightarrow c_{\mathrm{E}}\), this gives
\begin{equation}
    r_{\mathrm{cap}}
    \left(\frac{\pi}{2}\right)
    \rightarrow
    2R_{\mathrm{QG}} .
\end{equation}

This result is a TFT-native capacity surface. It should not be identified prematurely with the exact Kerr ergosurface. Rather, it is the mechanical surface on which torsional Soft-Floor drift equals the local propagation capacity. A completed Kerr correspondence must determine how this capacity surface maps onto the observational ergosphere after the horizon-scale normalization is fixed.

\section{Orbital and Photon-Ring Dynamics in the Rotating Exterior}

The eikonal propagation law supplies the starting point for photon-ring and shadow calculations. Define the massless rotating-exterior Hamiltonian
\begin{equation}
    \mathcal{H}_{\gamma}
    =
    \frac{1}{2}
    \left[
    -
    \frac{
    \left(
    \omega
    -
    m_{\phi}\omega_{\mathrm{drag}}(r)
    \right)^2
    }{
    c_P^2(r)
    }
    +
    k_r^2
    +
    \frac{k_{\theta}^2}{r^2}
    +
    \frac{m_{\phi}^2}{r^2\sin^2\theta}
    \right].
\end{equation}
Massless Soft-Floor shear rays satisfy
\begin{equation}
    \mathcal{H}_{\gamma}=0.
\end{equation}

In the equatorial plane, \(\theta=\pi/2\) and \(k_{\theta}=0\). The radial condition becomes
\begin{equation}
    k_r^2
    =
    \frac{
    \left(
    \omega
    -
    m_{\phi}\omega_{\mathrm{drag}}(r)
    \right)^2
    }{
    c_P^2(r)
    }
    -
    \frac{m_{\phi}^2}{r^2} .
\end{equation}
A circular photon-ring analogue occurs when the radial motion has a double root:
\begin{equation}
    k_r^2=0,
    \qquad
    \frac{d}{dr}
    \left(k_r^2\right)=0 .
\end{equation}
These two equations define the TFT photon-ring condition in the leading rotating exterior.

For massive boundary defects or compact orbiting matter, the same torsional field enters as a spin-orbit coupling. A minimal effective-potential form is
\begin{equation}
    \Phi_{\mathrm{eff}}
    =
    \Phi_{\mathrm{radial}}
    +
    \Phi_{\mathrm{orbital}}
    +
    \Phi_{\mathrm{spin}},
\end{equation}
with the leading spin term sourced by the Soft-Floor dragging rate:
\begin{equation}
    \Phi_{\mathrm{spin}}
    =
    -
    \ell_{\mathrm{orb}}
    \omega_{\mathrm{drag}}(r).
\end{equation}
Here \(\ell_{\mathrm{orb}}\) is the specific orbital angular momentum of the test body. The sign of this term distinguishes prograde and retrograde motion:
\begin{align}
    \ell_{\mathrm{orb}}\omega_{\mathrm{drag}}>0
    &\quad\Rightarrow\quad
    \text{prograde lowering of the effective strain barrier},
    \\
    \ell_{\mathrm{orb}}\omega_{\mathrm{drag}}<0
    &\quad\Rightarrow\quad
    \text{retrograde raising of the effective strain barrier}.
\end{align}

An ISCO analogue must then be derived from the ordinary stability conditions on the rotating effective potential:
\begin{equation}
    \frac{d\Phi_{\mathrm{eff}}}{dr}=0,
    \qquad
    \frac{d^2\Phi_{\mathrm{eff}}}{dr^2}=0 .
\end{equation}
The first equation selects a circular orbit, while the second marks the transition from stable to unstable circular motion. This appendix does not assert the exact Kerr ISCO formula. It supplies the TFT-native mechanical source of the prograde/retrograde split and identifies the equations from which the corresponding orbital radii must be derived.


\section{Rotating Phase-Boundary Geometry}
A rotating saturated boundary need not remain exactly spherical. Torsional loading modifies the stress distribution over the FCC/SC transition surface and can generate an axisymmetric deformation.

A minimal representation is
\begin{equation}
    R(\theta) = R_{\mathrm{QG}} + \delta R_{\mathrm{spin}}(\theta),
\end{equation}
where $R_{\mathrm{QG}}$ is the average saturation radius and $\delta R_{\mathrm{spin}}$ is the spin-induced boundary deformation.

Since the static boundary shape must be invariant under reversal of the spin direction, the leading deformation is even in $a_\ast$. The lowest-order axisymmetric correction is therefore expected to have the form
\begin{equation}
    \delta R_{\mathrm{spin}}(\theta) \propto a_\ast^2 P_2(\cos\theta),
\end{equation}
with $P_2$ the second Legendre polynomial.

This expression is an ansatz for the leading deformation channel. A Kerr-compatible TFT solution must determine whether saturation and relaxation select a unique spin-deformed boundary, rather than allowing arbitrary elastic surface structure.

\section{Multipole Constraint and Boundary Relaxation}

The strongest condition on a TFT Kerr correspondence limit is not the spin bound alone, but the exterior multipole structure.

A generic rotating continuum body can carry arbitrary shape-dependent multipoles. A Kerr black hole does not have freely adjustable exterior multipoles; its exterior field is fixed by mass and angular momentum. In TFT terms, the relaxed rotating FCC/SC boundary must therefore erase arbitrary surface memory and produce an exterior strain field determined only by \(M\) and \(\mathcal{J}_{\mathrm{BH}}\).

The required condition is
\begin{equation}
    \text{relaxed rotating FCC/SC boundary}
    \quad
    \Rightarrow
    \quad
    \text{exterior field fixed by } M \text{ and } \mathcal{J}_{\mathrm{BH}}.
\end{equation}

Equivalently, a Kerr-compatible saturated boundary must not retain independent surface hair in its relaxed exterior field.

This condition can be expressed as a dissipative boundary-selection principle. Let the stored exterior-plus-boundary strain energy be written schematically as
\begin{equation}
    \mathcal{E}_{\mathrm{stored}}
    =
    \mathcal{E}_{\mathrm{radial}}
    +
    \mathcal{E}_{\mathrm{tor}}
    +
    \mathcal{E}_{\partial},
\end{equation}
where \(\mathcal{E}_{\mathrm{radial}}\) is the static radial strain energy, \(\mathcal{E}_{\mathrm{tor}}\) is the exterior torsional Soft-Floor shear energy, and \(\mathcal{E}_{\partial}\) is the deformation energy of the live FCC/SC boundary.

During merger, accretion, or spin readjustment, arbitrary surface distortions may be excited. These distortions are not independent conserved exterior charges. They are transient boundary stresses that must be removed through stress shedding and ringdown. The dissipative approach to a relaxed exterior may be written as
\begin{equation}
    \frac{d\mathcal{E}_{\mathrm{stored}}}{dt}
    =
    -
    \mathcal{P}_{\mathrm{shed}}
    \leq
    0,
\end{equation}
where \(\mathcal{P}_{\mathrm{shed}}\geq0\) is the power carried away by outgoing gravitational stress waves and other boundary-relaxation channels.

The relaxed boundary satisfies
\begin{equation}
    \delta \mathcal{E}_{\mathrm{stored}}
    =
    0,
\end{equation}
subject to fixed total mass and fixed axial angular momentum,
\begin{equation}
    \delta M=0,
    \qquad
    \delta \mathcal{J}_{\mathrm{BH}}=0.
\end{equation}

Thus, non-universal surface multipoles are damped by stress shedding until the exterior field depends only on the conserved quantities \(M\) and \(\mathcal{J}_{\mathrm{BH}}\). This is the TFT boundary-relaxation analogue of the no-hair condition. Without such a selection principle, TFT would describe a generic rotating elastic object rather than a Kerr-correspondence black hole.

\section{Ringdown Spectrum of a Rotating Boundary}

The live-boundary appendix introduced a scalar relaxation model for post-merger boundary deformation. A Kerr-compatible limit requires a mode spectrum for the rotating boundary.

A general angular-mode expansion may be written as
\begin{equation}
    \delta R(\theta,\phi,t) = \sum_{\ell,m} A_{\ell m} e^{-\Gamma_{\ell m}^{(\mathrm{rd})}t} \cos \left( \omega_{\ell m}t+\varphi_{\ell m}^{(0)} \right) Y_{\ell m}(\theta,\phi).
\end{equation}

For a Kerr-like boundary attractor, the final ringdown spectrum must depend only on the remnant mass and spin:
\begin{align}
    \omega_{\ell m} &= \omega_{\ell m}(M,a_\ast), \\
    \Gamma_{\ell m}^{(\mathrm{rd})} &= \Gamma_{\ell m}^{(\mathrm{rd})}(M,a_\ast).
\end{align}

This is the dynamical counterpart of the multipole constraint. Merger may excite arbitrary transient surface structure, but the relaxed spectrum must lose dependence on uncontrolled boundary details. The observable ringdown modes then become normal modes of the rotating FCC/SC phase boundary.

\section{Horizon-Scale Normalization}

A Kerr-correspondence construction must distinguish the physical TFT phase boundary from the coordinate surfaces used in the standard Kerr description.

In the present TFT convention, the saturation radius is
\begin{equation}
    R_{\mathrm{QG}}
    =
    \frac{GM}{c_{\mathrm{E}}^2}.
\end{equation}

This is the physical FCC/SC phase boundary. It is the radius at which the active Soft Floor reaches the saturation condition
\begin{equation}
    \alpha_G
    =
    \frac{GM}{r c_{\mathrm{E}}^2}
    =
    1.
\end{equation}

Thus,
\begin{equation}
    R_{\mathrm{QG}}
    =
    \text{physical FCC/SC saturation boundary}.
\end{equation}

This surface should not be automatically identified with the Boyer--Lindquist coordinate horizon of Kerr geometry. In standard Kerr notation, the gravitational length is
\begin{equation}
    R_g
    =
    \frac{GM}{c^2},
\end{equation}
and the outer Kerr horizon coordinate is
\begin{equation}
    r_+
    =
    R_g
    \left(
    1+\sqrt{1-a_\ast^2}
    \right).
\end{equation}

For a nonrotating Schwarzschild black hole,
\begin{equation}
    r_+
    =
    2R_g,
\end{equation}
while for the extremal Kerr coordinate limit,
\begin{equation}
    r_+
    =
    R_g.
\end{equation}

TFT does not require the physical FCC/SC boundary to coincide directly with these coordinate radii. Instead, a Kerr correspondence requires two distinct notions of boundary:
\begin{equation}
    R_{\mathrm{QG}}
    =
    \text{physical phase-transition boundary},
\end{equation}
and
\begin{equation}
    R_{\mathrm{obs}}(a_\ast)
    =
    \text{outer optical or causal correspondence surface}.
\end{equation}

The physical boundary \(R_{\mathrm{QG}}\) is set by lattice saturation. The observational surface \(R_{\mathrm{obs}}(a_\ast)\) is the surface inferred from external light propagation, causal trapping, frame dragging, shadow structure, or accretion-disk phenomenology.

This distinction is necessary because TFT replaces the interior and horizon ontology of GR. The correspondence problem is therefore not
\begin{equation}
    R_{\mathrm{QG}}
    =
    r_+
\end{equation}
as an assumption. Rather, the required map is
\begin{equation}
    \left[
    R_{\mathrm{QG}},
    \,
    \omega_{\mathrm{drag}}(r),
    \,
    c_P(r),
    \,
    v_{\phi}^{(\mathrm{SF})}(r,\theta)
    \right]
    \longrightarrow
    R_{\mathrm{obs}}(a_\ast).
\end{equation}

Thus, the completed rotating TFT black-hole solution must determine how the physical FCC/SC boundary and the derived optical/capacity surfaces correspond to the observational horizon, ergosphere, photon ring, and shadow scales used in Kerr phenomenology.

\section{Minimal Kerr-Correspondence Program}

A TFT continuum construction of a Kerr-correspondence exterior requires the following minimum steps. The objective is not to reproduce Kerr geometry as a fundamental ontology, but to determine when the rotating TFT boundary produces the same successful external observational signatures.

First, specify the rotating black-hole state by its conserved mass and macroscopic axial angular momentum,
\begin{equation}
    M, \qquad \mathcal{J}_{\mathrm{BH}},
\end{equation}
with dimensionless spin parameter
\begin{equation}
    a_\ast = \frac{c_{\mathrm{E}}\mathcal{J}_{\mathrm{BH}}}{GM^2}.
\end{equation}

Second, use the live-boundary spin normalization
\begin{equation}
    \mathcal{J}_{\mathrm{BH}} = I_{\mathrm{surf}} \frac{v_{\mathrm{tor}}}{R_{\mathrm{QG}}},
\end{equation}
so that the observed spin is carried mechanically by torsional loading of the FCC/SC phase boundary and by shear strain in the surrounding Soft Floor.

Third, use the derived axisymmetric torsional Soft-Floor field equation
\begin{equation}
    \frac{\mathrm{d}}{\mathrm{d}r} \left[ \tau_{\mathrm{tor}}(r) r^4 \frac{\mathrm{d}\omega_{\mathrm{drag}}}{\mathrm{d}r} \right] = 0
\end{equation}
to determine the exterior angular-drag field. In the constant-stiffness exterior limit this gives
\begin{equation}
    \omega_{\mathrm{drag}}(r) = \omega_{\partial} \left( \frac{R_{\mathrm{QG}}}{r} \right)^3,
\end{equation}
where
\begin{equation}
    \omega_{\partial} = \frac{v_{\mathrm{tor}}}{R_{\mathrm{QG}}}.
\end{equation}

Fourth, connect the dragging field to wave propagation through the rotating optical relation
\begin{equation}
    \left[ \omega - m_{\phi}\omega_{\mathrm{drag}}(r) \right]^2 = c_P^2(r) \left[ k_r^2 + \frac{k_{\theta}^2}{r^2} + \frac{m_{\phi}^2}{r^2\sin^2\theta} \right].
\end{equation}
This equation supplies the TFT frame-dragging mechanism: prograde and retrograde modes experience different effective propagation conditions because the local Soft Floor is torsionally dragged.

Fifth, identify the torsional capacity surface by solving
\begin{equation}
    \left| v_{\phi}^{(\mathrm{SF})}(r,\theta) \right| = c_P(r),
\end{equation}
rather than postulating an ergosphere by analogy. In the constant-stiffness exterior approximation,
\begin{equation}
    v_{\phi}^{(\mathrm{SF})}(r,\theta) = v_{\mathrm{tor}} \left( \frac{R_{\mathrm{QG}}}{r} \right)^2 \sin\theta,
\end{equation}
and
\begin{equation}
    c_P(r) = c_{\mathrm{E}} \left( 1 - \frac{R_{\mathrm{QG}}}{r} \right)^2.
\end{equation}
Therefore, the capacity surface is
\begin{equation}
    r_{\mathrm{cap}}(\theta) = R_{\mathrm{QG}} \left[ 1 + \sqrt{ \frac{v_{\mathrm{tor}}}{c_{\mathrm{E}}} \sin\theta } \right].
\end{equation}

Sixth, determine the spin-induced boundary deformation by solving the live-boundary equilibrium problem. At lowest order, the deformation is expected to enter through an axisymmetric quadrupolar channel,
\begin{equation}
    R(\theta) = R_{\mathrm{QG}} + \delta R_{\mathrm{spin}}(\theta),
\end{equation}
with
\begin{equation}
    \delta R_{\mathrm{spin}}(\theta) \propto a_\ast^2 P_2(\cos\theta)
\end{equation}
as the leading deformation ansatz. A completed theory must derive this deformation from boundary stress balance rather than assume it.

Seventh, impose the boundary-relaxation selection principle. Arbitrary surface multipoles must be damped by stress shedding and ringdown, leaving a relaxed exterior determined only by
\begin{equation}
    M \qquad \text{and} \qquad \mathcal{J}_{\mathrm{BH}}.
\end{equation}
Equivalently, the relaxed boundary must satisfy
\begin{equation}
    \delta \mathcal{E}_{\mathrm{stored}} = 0,
\end{equation}
with dissipative approach
\begin{equation}
    \frac{\mathrm{d}\mathcal{E}_{\mathrm{stored}}}{\mathrm{d}t} = - \mathcal{P}_{\mathrm{shed}} \leq 0.
\end{equation}

Eighth, derive ISCO, photon-ring, and shadow analogues from the same rotating exterior propagation law. For massless shear excitations, the eikonal Hamiltonian is
\begin{equation}
    \mathcal{H}_{\gamma} = \frac{1}{2} \left[ - \frac{ \left( \omega - m_{\phi}\omega_{\mathrm{drag}}(r) \right)^2 }{ c_P^2(r) } + k_r^2 + \frac{k_{\theta}^2}{r^2} + \frac{m_{\phi}^2}{r^2\sin^2\theta} \right] = 0.
\end{equation}
Photon-ring and shadow structures must follow from stationary and unstable circular solutions of this Hamiltonian, rather than from an imported Kerr metric.

Ninth, derive the rotating-boundary ringdown spectrum as a normal-mode problem of the live FCC/SC boundary coupled to the torsional Soft-Floor exterior:
\begin{equation}
    \omega_{\ell m} = \omega_{\ell m}(M,a_\ast), \qquad \Gamma_{\ell m}^{(\mathrm{rd})} = \Gamma_{\ell m}^{(\mathrm{rd})}(M,a_\ast).
\end{equation}
A Kerr-correspondence limit requires that the dominant relaxed modes depend only on the conserved mass and spin, not on arbitrary initial boundary distortions.

Tenth, determine the correspondence map between the physical TFT phase boundary and the observational horizon-scale surfaces:
\begin{equation}
    R_{\mathrm{QG}} \longrightarrow R_{\mathrm{obs}}(a_\ast).
\end{equation}
Here
\begin{equation}
    R_{\mathrm{QG}} = \frac{GM}{c_{\mathrm{E}}^2}
\end{equation}
is the physical FCC/SC saturation boundary, while $R_{\mathrm{obs}}(a_\ast)$ denotes the outer optical or causal surface inferred from light propagation, frame dragging, capacity trapping, shadow structure, or accretion phenomenology.

These conditions define the minimum mathematical content of a Kerr-correspondence TFT solution. The completed theory must show that the physical FCC/SC phase boundary, the torsional Soft-Floor drag field, and the rotating optical propagation law reproduce the empirically successful external signatures of Kerr black holes while replacing the GR interior and horizon ontology with mechanical phase-boundary dynamics.

\section{Conclusion}
This appendix formulates the Kerr-limit problem for TFT continuum mechanics. The live-boundary model provides the mechanical carrier of black-hole angular momentum. A Kerr-compatible extension must derive the exterior rotating field produced by that carrier.

The proposed structure is
\begin{equation}
    \text{rotating black hole} = \text{phase-locked saturated domain} + \text{torsional FCC/SC boundary} + \text{external rotating Soft-Floor shear field}.
\end{equation}

The corresponding development sequence is
\begin{equation}
    \text{boundary torsion} \rightarrow \text{Soft-Floor shear} \rightarrow \text{rotating optical geometry} \rightarrow \text{capacity surface} \rightarrow \text{ISCO/photon-ring dynamics} \rightarrow \text{ringdown}.
\end{equation}

The present appendix identifies the Kerr-correspondence conditions required of a rotating TFT exterior. The detailed constitutive closure, coefficient matching, photon-ring equations, ISCO map, observational surface map, ringdown scaling, and numerical exterior solver are developed in \hyperref[app:bh_rotating_kerr_clouser]{Appendix~\ref*{app:bh_rotating_kerr_clouser}}.